\section{Threshold Signatures}
\label{threshold-sig-background}

This chapter defines threshold signature schemes.
Portions of this chapter are adapted from previous work by Crites et al.~\cite{CritesKM2023}.

To simplify notation, the syntax assumes a centralized key generation algorithm to generate the joint public key representing all $n$ signers and the shares of the corresponding secret key.
However, a distributed key generation (DKG) protocol may replace this trusted dealer.


\begin{definition} \label{def:threshold-sig}
  A \defn{threshold signature scheme} $\threshsig$ whose signing protocol consists of \numrounds~rounds is a tuple of the following algorithms
  $\threshsig = (\Setup, \KeyGen, \allowbreak (\Sign_1, \ldots, \Sign_\numrounds ), \allowbreak \Combine, \allowbreak \Verify, \Recover)$, defined as follows:


\begin{itemize}

  \item $\Setup(\usecp) \rightarrow \pp$: Takes as input a security parameter and outputs public parameters $\pp$ (which are given implicitly as input to all other algorithms).

  \item $\KeyGen(n,\thresh) \rightarrow (\aggpk, \{ \aggpk_i \}_{i \in \set{n}}, \{ \sk_i \}_{i \in \set{n}})$:
    A probabilistic algorithm that takes as input the number of signers $n$ and the threshold $\thresh$ and outputs the public key $\aggpk$ representing the set of $n$
signers and the set $\{ \aggpk_i \}_{i \in \set{n}}$ of public keys for each signer,
    such that $\aggpk_i \in \pubkeydomain$ and $\aggpk \in \pubkeydomain$,
    where $\pubkeydomain$ is the domain of public keys for $\threshsig$.
    The algorithm also outputs set $\{ \sk_i \}_{i \in \set{n}}$ of secret shares corresponding to a secret key $\sk$,
    where $\sk_i \in \seckeydomain$ and $\sk \in \seckeydomain$,
    such that $\seckeydomain$ is the domain of secret keys for $\threshsig$.
    Each $\sk_i$ is given to participant $i$, for $i \in \set{n}$.
    It is assumed that participant $i$ is sent its respective share $\sk_i$ privately.

  \item $(\Sign_1, \ldots, \Sign_\numrounds) \rightarrow \{\msg_{ 1, i},
    \ldots, \msg_{\numrounds, i} \}_{i \in \coalition}$:
    A set of probabilistic algorithms where each subsequent algorithm represents a single stage in an interactive signing protocol,
    performed by each signing party in a signing set
    $\coalition \subseteq \set{n}, |\coalition| \geq \thresh$,
    with respect to a message $m$, defined as follows:
  \begin{align}\label{eqn:protocol-messages}
    \begin{split}
    (\msg_{1,i}, \state_{1,i}) &\gets \Sign_1(i, \sk_i, \coalition, m)   \\
    (\msg_{2,i}, \state_{2,i}) &\gets \Sign_2(\state_{1,i}, \PM_{1})   \\
    \vdots \\
    \msg_{\numrounds, i} &\gets \Sign_\numrounds(\state_{\numrounds-1,i}, \PM_{\numrounds-1}) \\
    \end{split}
  \end{align}
    where $\msg_{ 1, i}, \ldots, \msg_{\numrounds, i}$ are protocol messages for each signing round,
     $\PM$ is a set of protocol messages, and
  $\state_{j,i}$ is the state of signing party $i$ in $\coalition$ in round $j$.
  The signing set $\coalition$ and message $m$ are shown here as inputs to the first round,
  but they may be deferred to subsequent rounds,
  depending on the scheme.

\item $\Combine(\coalition, m, \PM_1, \PM_2, \dots, \PM_\numrounds) \rightarrow (m, \sigma)$:
  A deterministic algorithm that takes as input the signing set $\coalition$, the message $m$, and a set of protocol messages
  $\PM_1, \PM_2, \dots, \PM_\numrounds$ from all $\numrounds$ signing rounds
  and outputs a signature $\sigma$. (Note that this algorithm is separate from the signing rounds and may be performed by one of the signers or an external party.)

  \item $\Verify(\threshpk, m, \sigma) \rightarrow 0/1$: A deterministic algorithm that takes as input the public key $\threshpk$, a message $m$, and a purported signature $\sigma$ and outputs $1$ (accept) if the signature verifies; else, it outputs $0$ (reject).

  \item $\Recover(\thresh, \{ (i, \sk_i)\}_{ i \in \coalition}) \rightarrow \bot / \sk$:
    A deterministic algorithm that takes as input a threshold $\thresh$ and a set of secret key shares $\{(i, \sk_i) \}_{i \in \coalition}$.
    If $\coalition \not\subseteq \set{n}$ or $\lvert \coalition \rvert < \thresh$,
    it outputs $\bot$; else, it recovers $\sk$ using the set of shares.
\end{itemize}
\end{definition}

\begin{remark}
The syntax assumes that the signing set $\coalition$, message $m$, and secret key share $\sk_i$ are given as input in the first round of signing.
  However, constructions may defer these values to later rounds of signing,
  if the proof of security so allows.
\end{remark}

\begin{remark}[On the Inclusion of an Explicit Recovery Algorithm]
  We define a threshold signature scheme with an explicit recovery algorithm.
  While this recovery algorithm is defined explicitly for centralized and decentralized threshold secret sharing schemes~\cite{BonehS2020,GennaroJKR07},
  its use for threshold signatures is often implicit.
  We define this recovery algorithm explicitly,
  and later employ it to prove our results.
\end{remark}

\paragraph{Correctness.}
A threshold signature scheme is \emph{correct} if for all security parameters $\secp$, all allowable $1 \leq \thresh \leq n$,
all $\coalition \subseteq [n]$ such that $| \coalition | \geq \thresh$,
all messages $m$ and
for $(\aggpk, \{ \aggpk_i \}_{i \in \set{n}}, \{ \sk_i \}_{i \in \set{n}}) \allowbreak \randpick \KeyGen(n, \allowbreak \thresh)$,
the $\Recover$ algorithm outputs $\sk$ and the experiment in (\ref{eqn:protocol-messages}) returns signature shares that, when combined with all other shares, produce a signature $\sigma$
such that $\Verify(\pk, m, \sigma) = 1$.

\subsection{Static Security}

\begin{figure}[!htb]
	\begin{pchstack}[boxed]
		\begin{pcvstack}
      \procedure{$\TUnforgExp_{\A,\threshsig}(\secp, n, \thresh, \tcornew)$}{
		%\protect\dashbox[3.5cm][c]{$\G^{\aduf}_{\A,\threshsig}(\secp, \hmulp)$}}{
			\sessions \gets \emptyset \\
			\pp \leftarrow  \ParGen(\usecp) \\
      			%Q \gets \emptyset \pccomment{set of queried messages} \\ Liz: Do all or none?
 	    		(\corrupt, \state_\A) \randpick \A(\pp, n, \thresh, \tcornew)   \\
			\pcreturn \bot~\pcif \corrupt \not\subset [n] \vee |\corrupt| > \tcornew \\
			%\text{\algcom or \dbox{$\vert \corrupt \rvert \geq t/2$ (Theorem~\ref{thm:OMDL+ROM})} } \\
        			\honest \gets \set{n} \setminus \corrupt  \\
			%\state_{k,\ssid} \gets ( )  \text{\algcom state vectors for honest signers} \\
			% CK- these are created as the output from Sign
       			(\threshpk, \{ \aggpk_i, \sk_i \}_{i \in \set{n}}) \randpick  \KeyGen(n,\thresh) \\
        			\mathsf{input} \gets (\threshpk, \{ \aggpk_i \}_{i \in \set{n}},  \{\sk_j \}_{j \in \corrupt}, \state_\A) \\
              		%\fbox{ $(m^*, \sig^*) \randpick \A^{\BigO^{\Sign, \Sign', \Sign''}}(\mathsf{input}) $}\\
              		(m^*, \sig^*) \randpick \A^{\BigO^{\tiny{\Sign_1, \Sign_2, \dots, \Sign_\numrounds, \fbox{$\Corrupt$}}}}(\mathsf{input}) \\
      \pcreturn 1~\pcif m^* \notin Q_m \\
      ~~~\wedge  \Verify(\threshpk, m^*, \sig^*) = 1\\
			\pcreturn 0
		}
   		\pcvspace
		\procedure{$\Init(\ssid, k, \coalition, m)$}{
		\pcif \ssid \notin \sessions \\
		~~~\pcreturn \bot~\pcif \coalition \not\subseteq [n] \vee |\coalition| \neq \thresh\\
		~~~~\vee k \notin \honest \cap \coalition \\
		~~~\sessions \gets \sessions \cup \{\ssid\} \\
		~~~\coalition \gets \coalition[\ssid];~m \gets m[\ssid] \\
		~~~Q_m \gets Q_m \cup \{m[\ssid]\} \\
		\pcreturn \bot~\pcif k \notin \honest \cap \coalition[\ssid] \\
		~~~\vee \round[\ssid, k] \neq \bot
		}
		\end{pcvstack}
    		\pchspace
		\begin{pcvstack}
			\procedure{$\BigO^{\Sign_1}(\ssid, k, \coalition, m)$}{
       	 			%\pccomment{$k$ denotes an honest participant identifier} \\
				\pcreturn \bot~\pcif \bot \gets \Init(\ssid, k, \coalition, m) \\
				%\S \gets \signers[\ssid];~m \gets m[\ssid] \\
				(\PM_1[\ssid, k], \state_1[\ssid, k]) \\
				~~~~~~~~~~~\gets  \Sign_1(k, \sk_k, \coalition[\ssid], m[\ssid]) \\
				%\PM_1[\ssid, k] \gets \PM_1;~\state_1[\ssid, k] \gets \state_1 \\
				\PM_{1, k} \gets \PM_1[\ssid, k];~\round[\ssid, k] \gets 1 \\
        \pcreturn \PM_{1, k}
			}
    		\pcvspace
    		\procedure{$\BigO^{\Sign_2}(\ssid, k, \PM_1)$}{
			\pcreturn \bot~\pcif \ssid \notin \sessions \\
			~~~\vee k \notin \honest \cap \coalition[\ssid] \vee \round[\ssid, k] \neq 1 \\
			\pcparse \{\PM_{1,i}\}_{i \in \coalition} \gets \PM_1 \\
			\pcreturn \bot~\pcif \PM_{1, k} \neq \PM_1[\ssid, k] \\
			%\state_1 \gets \state_1[\ssid, k] \\
				(\PM_2[\ssid, k], \state_2[\ssid, k]) \\
				~~~~~~~~~~~\gets  \Sign_2(\state_1[\ssid, k], \PM_1) \\
				%\PM_2[\ssid, k] \gets \PM_2;~\state_2[\ssid, k] \gets \state_2 \\
				\PM_{2, k} \gets \PM_2[\ssid, k];~\round[\ssid, k] \gets 2 \\
       		 	\pcreturn \PM_{2,k}
		}
		\vdots
		\procedure{$\BigO^{\Sign_\numrounds}(\ssid, k, \PM_{\numrounds-1})$}{
			\pcreturn \bot~\pcif \ssid \notin \sessions \\
			~~~\vee k \notin \honest \cap \coalition[\ssid] \vee \round[\ssid, k] \neq \numrounds-1 \\
			\pcparse \{\PM_{\numrounds-1,i}\}_{i \in \coalition} \gets \PM_{\numrounds-1} \\
			\pcreturn \bot~\pcif \PM_{\numrounds-1, k} \neq \PM_{\numrounds-1}[\ssid, k] \\
			%\state_2 \gets \state_2[\ssid, k] \\
				\PM_\numrounds[\ssid, k] \gets  \Sign_\numrounds(\state_{\numrounds-1}[\ssid, k], \PM_{\numrounds-1}) \\
				%\PM_3[\ssid, k] \gets \PM_3;~
				\PM_{\numrounds, k} \gets \PM_\numrounds[\ssid, k];~\round[\ssid, k] \gets \numrounds \\
       		 	\pcreturn \PM_{\numrounds, k}
		}
		\end{pcvstack}
	\end{pchstack}
\caption{Static and adaptive unforgeability games for a threshold signature scheme with $\numrounds$ signing rounds.
}
\label{fig:euf-cma-threshold}
\hrulefill
\end{figure}



Most threshold signature schemes in the literature are proven in the \emph{static} security model,
which captures a limited adversary that can only corrupt parties at the onset of the protocol.
We provide a game-based definition in \ref{fig:euf-cma-threshold}, inspired by the definitions in \cite{CritesKM23,LibertJY16}.
The static security game contains all but the boxes and is analogous to existential unforgeability under chosen message attack (EUF-CMA)~\cite{GoldwasserMR88} for standard, single-party signature schemes;
adaptive game adds the boxes.
The public parameters $\pp$ are implicitly given as input to all algorithms, and $\PM_1, \PM_2, \dots, \PM_\numrounds$ represent protocol messages defined within the construction.

In the static unforgeability game,
the environment accepts a security parameter $\secp$,
total number of players $n$, a reconstruction threshold $\tcor$,
and a corruption threshold $\tcornew$ as input,
such that $\tcornew < t \leq n$.
%The number of parties that the adversary corrupts must stay within the bound $\hmulp \cdot \tcor$.
%Typically, a static adversary corrupts a full $\tcor$ (i.e., $\hmulp = 1$) signers out of at least $\thresh$ signers in a session.
$\tcornew$ specifies the number of signers that the adversary may corrupt,
and $\tcor$ the reconstruction threshold,
such that $\tcornew \leq \tcor \leq n$.

The environment generates public parameters $\pp$
and initializes the adversary with $(\pp, n, \tcor)$.
The adversary chooses a set of corrupted parties $\corrupt$,
and the challenger sets the honest parties $\honest \gets [n] \setminus \corrupt$.
The challenger then runs $\KeyGen$ to derive the joint public key $\aggpk$ representing all $n$ signers, the individual public key shares $\{ \aggpk_i \}_{i \in \set{n}}$, and the secret key shares $ \{ \sk_i \}_{i \in \set{n}}$.
It returns $\aggpk$, $\{ \aggpk_i \}_{i \in \set{n}}$, and the set of corrupt signing shares $ \{ \sk_j \}_{j \in \corrupt}$ to the adversary.


After key generation has concluded,
the adversary can then query signing oracles $\BigO^{\Sign_1}, \ldots, \BigO^{\Sign_\numrounds}$ for honest signers to engage in subsequent signing rounds, across different sessions.
The environment performs various checks during the signing rounds to ensure each participant being called is in fact uncorrupted,
is in the correct signing round,
and that their protocol message in each round is faithfully carried to the next.
Each of these checks is carried out with respect to the signing session identifier, denoted by $\ssid$.
%
The adversary wins if it can produce a valid forgery $\sigma^* = (R^*, z^*)$ with
respect to the joint public key $\aggpk$ on a message $m^*$ that has not been previously queried.
%
We model the adversary as \emph{rushing},
as it may wait to produce its outputs after having seen honest protocol messages in each round.
The adversary is allowed to be \emph{concurrent},
as it may open simultaneous signing sessions at once
or choose not to complete a signing session.
We model concurrency explicitly in our definition to ensure that our notion of security protects against practical attacks which can occur under a concurrent adversary~\cite{BenhamoudaLLOR20}.

\begin{definition}[Static Security]
Let the advantage of an adversary $\A$ playing the static security game
$\G^{\uf}_{\A,\threshsig}(\secp, n, \tcor, \hmulp)$, as defined in Figure~\ref{fig:euf-cma-threshold}, be as follows:
%
\[
  \advant{$\uf$}{\A,\threshsig}(\secp, n, \tcor, \tcornew) = \bigabs{\Pr[\G^{\uf}_{\A,\threshsig}(\secp, n, \tcor, \tcornew) = 1]}
\]

A threshold signature scheme $\threshsig$ is \emph{statically secure} if for all PPT adversaries $\A$,
$\advant{$\uf$}{\A}(\secp, n, \tcor,\tcornew)$ is negligible,
for $(n, \tcor, \tcornew) \in \mathbb{N}, \tcor \leq n$, $\tcornew < \tcor$.
\end{definition}


\subsection{Adaptive Security}

In this work,
we consider only static security for our threshold Schnorr signature constructions.
We provide a high-level discussion of adaptive security here,
for comparison.


%As in static security, the adversary can corrupt a \emph{full} $\tcor$ out of $\thresh$ parties, and it is assumed that there is at least one honest signer.
In the adaptive setting, the adversary is given the added capability to corrupt parties \emph{at any point during the protocol}, based on all protocol messages it has seen.
This is in contrast to the static setting, where the adversary must declare the set of corrupt parties at the onset of the protocol, before any messages have been sent.
Adaptive corruption means the adversary receives not only the honest party's secret key, but also its \emph{state across all signing sessions.}

At any time throughout the experiment, the adversary may query $\BigO^{\Corrupt}$ to corrupt
up to $\tcornew$ honest signers,
obtaining \emph{both} their secret key and as well as state across all signing sessions.
As such,
in the adaptive setting,
the adversary is not restricted to choosing its set of
corrupt participants $\corrupt$ only at the beginning of the game,
and learns additional information throughout the game (e.g., the signing state for each corrupted participant).




